\documentclass[conference]{IEEEtran}
\usepackage[T1]{fontenc}
\usepackage{booktabs}
\usepackage{amsmath}
\usepackage{url}
\usepackage{xcolor}
\usepackage[hidelinks]{hyperref}

% Marks a number that still has to come from the final benchmark run.
\newcommand{\tbd}[1]{\textcolor{red}{[#1]}}

\begin{document}

\title{What Is a Learned Index Worth Inside an Embedded SQL Engine?\\
A Controlled Comparison with a B+-Tree}

\author{\IEEEauthorblockN{Aadit Sood}
\IEEEauthorblockA{University of Maryland, College Park}}

\maketitle

\begin{abstract}
Learned indexes answer point lookups faster than B-trees when measured on
their own, but a database query also has to be parsed, planned, logged
and turned back into a row. We built a small embedded SQL engine whose
storage engine is a learned index (piecewise-linear segments over a
gapped array, with packed-memory-array rebalancing), and made the index
swappable at compile time for a well-known B+-tree library with nothing
else changed, write-ahead log and crash recovery included. On its own,
the learned index is 1.5--1.6x faster than the B+-tree at 1M keys and
1.1--1.4x at 10M, and between 0.91x and 1.30x on four real datasets.
Inside the engine the advantage shrinks to 1.12--1.18x for a query that
was parsed ahead of time and to 1.05--1.08x for one that is parsed each
time. Range scans and random-order inserts are even; appends are 2--3x
slower on the learned index, and inserts that do not follow the key
distribution are up to 3.4x slower. Durable writes are dominated by
fsync whichever index is used, even with batching; opening the database
from a snapshot is 1.4--1.6x faster on the learned index, and memory per
row is about the same. We also describe two performance problems that
stayed hidden while the index was compared only with weak baselines, and
the differential testing against SQLite that both builds pass.
\end{abstract}

\section{Introduction}

A learned index replaces the traversal of a B-tree with a model that
predicts where a key sits in a sorted array, followed by a short local
search to correct the prediction~\cite{kraska2018}. Measured on their own,
learned indexes answer lookups faster than B-trees and need far less memory
for their routing structure~\cite{marcus2020}. A database does much more
than look up keys, though. It parses and plans each statement, logs each
write before applying it, copies rows in and out of its storage, and has to
recover after a crash. How much of the index's advantage is left once all of
that is added is a different question from how fast the index is.

Several systems have tried to answer it for their own setting. Bourbon put
learned models into the files of an LSM-tree key-value store~\cite{dai2020},
Google measured learned indexes in a disk-based database~\cite{abulibdeh2020},
and PostLearn integrated the ALEX+ learned index into PostgreSQL as an index
access method~\cite{fuad2026}. PostLearn is the closest to this paper. It
reports speedups of up to 1.5x for point lookups inside PostgreSQL, smaller
ones once the client is included, and it supports reads only; updates,
write-ahead logging and crash recovery are left as future work.

This paper studies those missing parts in a much smaller setting: an
embedded SQL engine in the style of SQLite, built from scratch, whose
storage engine is a learned index. The engine has a recursive-descent SQL
parser, an executor that routes each query to a point lookup, range scan,
secondary-index lookup or full scan, a catalog stored as a table, a
checksummed write-ahead log, and snapshot checkpoints. It is built so that
the learned index can be replaced, at compile time, by a well-known B+-tree
library (tlx~\cite{tlx}) with nothing else changing. Both builds pass the
same test suites and the same differential test against
SQLite~\cite{sqlite}. That makes it possible to measure what the learned
index is worth at each layer of the engine and for each kind of operation:
point lookups with and without parsing, range scans, inserts in random and
increasing order, inserts that follow or depart from the existing key
distribution, durable writes with and without batched fsync, recovery time,
and memory per row.

The short answer is that most of the advantage does not survive. The
learned index is 1.5--1.6x faster than tlx at the storage layer on 1M
keys, but a point query through SQL is only 5--8\% faster, because the
lookup is a small part of what the query costs. Range scans and random
inserts are even. The learned index's known weaknesses, on the other
hand, carry through in full: appends are 2--3x slower, and inserts that
land where the data is sparse are up to 3.4x slower. Durable writes are
dominated by fsync on either index. The learned index's clearest gains
inside the engine are in startup time and in read-heavy workloads after
inserts, where the B+-tree's split leaves slow it down.

None of the index's techniques are new, and the paper does not claim them.
The segmentation is the greedy method of FITing-tree~\cite{galakatos2019},
the gapped array comes from ALEX~\cite{ding2020}, and the rebalancing scheme
is that of a packed-memory array~\cite{itai1981,bender2000}. The
contributions are:
\begin{itemize}
\item an embedded SQL engine whose index can be swapped between a learned
index and a B+-tree with everything else held fixed, with the write-ahead
log, checkpointing and recovery included (Section~\ref{sec:system});
\item a measurement of the learned index's value inside that engine across
reads, writes, durability, recovery and memory, on synthetic data and on
four real datasets from the SOSD benchmark (Section~\ref{sec:results});
\item an account of two performance problems that stayed invisible while
the index was compared only against weak baselines, and of the testing that
caught correctness bugs, offered as methodology for others building similar
systems (Section~\ref{sec:lessons}).
\end{itemize}
The code, the benchmark scripts and the raw output of every run reported
here are public.

\section{Background and Related Work}

\paragraph{Learned indexes} Kraska et al.\ proposed replacing index
structures with models that learn the cumulative distribution of the
keys~\cite{kraska2018}. FITing-tree splits the keys into linear segments
with a guaranteed maximum error, fitting them greedily in one
pass~\cite{galakatos2019}. The PGM-index computes the optimal segmentation
for a given error bound with O'Rourke's algorithm~\cite{orourke1981} and
indexes the segments recursively~\cite{ferragina2020}. RadixSpline builds a
spline and a radix table in a single pass~\cite{kipf2020}. ALEX made learned
indexes updatable by storing keys in gapped arrays, with room left for
inserts, and adapting its node layout to the data~\cite{ding2020}. SOSD
compares these and other structures on real datasets~\cite{marcus2020}.

\paragraph{Evaluations of updatable learned indexes} Wongkham et al.\
evaluated updatable learned indexes under many workloads and found their
advantages depend heavily on the workload and the data~\cite{wongkham2022}.
Lan et al.\ found that once data lives on disk, a B+-tree is competitive
with them in most settings~\cite{lan2023}.

\paragraph{Learned indexes in systems} Bourbon reported 1.23x to 1.78x
faster lookups in an LSM-tree store~\cite{dai2020}. Abu-Libdeh et al.\
describe learned indexes in Bigtable~\cite{abulibdeh2020}. PostLearn, as
described above, measured ALEX+ inside PostgreSQL for reads~\cite{fuad2026}.

\paragraph{Robustness} Learned indexes can be attacked through the data
they are built on. Kornaropoulos et al.\ poisoned the training data of
recursive model indexes~\cite{kornaropoulos2022}; Bachfischer et al.\
measured slowdowns of up to 20\% from poisoning in ALEX and the
PGM-index~\cite{bachfischer2022}; and Yang et al.\ built insert sequences
that slow ALEX down by up to 1,641x or exhaust its memory~\cite{yang2024}.

\paragraph{Testing and measuring} Differential testing of SQL engines goes
back to RAGS~\cite{slutz1998}; SQLancer applies related ideas at
scale~\cite{rigger2020}. Mytkowicz et al.\ showed how innocuous-looking
setup changes bias performance measurements~\cite{mytkowicz2009}, which is
the reason for the interleaved measurements in Section~\ref{sec:method}.

\section{The Engine}
\label{sec:system}

\subsection{Overview}

The engine is single-process and keeps all rows in memory, like an
in-memory SQLite database with a log. A statement is tokenized and parsed
by a recursive-descent parser; the executor turns an equality condition on
the primary key into a point lookup, \texttt{BETWEEN} on the primary key
into a range scan, an equality condition on an indexed column into a
secondary-index lookup, and anything else into a full scan with filtering.
The storage layer maps 64-bit keys to byte strings. Each row's key combines
a 16-bit table id with its 48-bit primary key, so a table's rows are
contiguous and ordered; secondary-index entries combine a 20-bit value with
a 28-bit primary key and live in the same store. The catalog is a table
(id 0) and rebuilds itself on startup by scanning it.

Every write is appended to the write-ahead log, with a CRC-32 per record,
before it is applied. Opening the log truncates it to its last intact
record. A checkpoint writes the whole store to a snapshot (to a temporary
file, fsync'd and renamed) and then empties the log. On startup the engine
loads the snapshot, folds the log into it (for each key, the last record
logged wins), and builds the index once from the result.

\subsection{The learned index}

The index stores keys in a gapped array that starts 70\% full. A model of
linear segments predicts each key's slot. Segments are fit in one pass:
each segment's line is pinned to its first key, and the segment grows while
some slope keeps every key within $\epsilon = 64$ slots of its prediction.
This is the greedy method of FITing-tree~\cite{galakatos2019}; an optimal
segmentation~\cite{orourke1981,ferragina2020} is also implemented but not
used, because the index refits regions often and the optimal fit is four to
seven times slower to compute. Segments are capped at 4,096 keys, because on
perfectly regular keys a single line would fit the whole array and
refreshing that one segment would mean rebuilding everything.

A lookup finds the segment with a branchless binary search over the
segments' first keys, computes the prediction, clamps it to the segment's
stretch of the array, and walks from there toward the first key that is not
smaller than the target, skipping empty slots. The walk takes at most
$\epsilon$ steps in each direction before switching to an exponential
search, so the worst case is logarithmic. Correctness never depends on the
model: because the array is sorted, the walk finds the right slot from any
starting point, and the model only decides how long the walk is.

\subsection{Inserts}

An insert shifts keys toward the nearest empty slot, which moves each
shifted key one slot away from where the model predicts it. Each segment
counts the inserts that changed it, meaning the segment that received the
key and any segment whose keys were shifted, and is refit once its count
reaches $\epsilon$. Since each counted insert moves a key by at most one
slot, every key the model was fit on stays within $2\epsilon$ of its
prediction. When a segment's region is too full to refresh (90\%), the
refresh widens to aligned windows of two, four, eight segments, with
density limits that fall from 0.9 for one segment to 0.8 for the whole
array, as in a packed-memory array~\cite{itai1981,bender2000}. Only an array
more than 80\% full is rebuilt, larger. The last segment can instead grow
the array at its end, by a quarter of its size at a time, which is what
keys arriving past the end use.

\subsection{Batched fsync}

By default every log record is fsync'd before the write returns. With a
setting $k > 1$, only every $k$-th record is fsync'd, and the log is synced
on close and before a checkpoint empties it. Every record still reaches the
operating system before the write returns, so a crash of the process loses
nothing; up to $k-1$ acknowledged writes can be lost if the machine loses
power. SQLite makes the same trade with \texttt{synchronous=NORMAL}. This
is not group commit in the usual sense, in which concurrent clients share
an fsync and each still waits for it: the engine has one client and no
multi-statement transactions.

\subsection{The B+-tree build}
\label{sec:btree}

The storage layer refers to its index through a single type alias. A
compile-time switch replaces the gapped array with a thin wrapper around
tlx's \texttt{btree\_map}~\cite{tlx}, which stores each row next to its key
in the tree's leaves, as the gapped array does in a parallel array. tlx's
B+-tree is the successor of the stx B+-tree used by SOSD. The parser,
executor, catalog, log and snapshot code are identical in both builds; the
only other difference is the text of \texttt{EXPLAIN}. Both builds pass the
same unit and end-to-end tests and agree with SQLite on the same randomized
differential test.

\section{Methodology}
\label{sec:method}

\paragraph{Machine} All measurements come from one laptop: an Intel Core
i7-13620H (6 performance and 4 efficiency cores, 24 MB of L3 cache) running
Linux 6.18 under WSL2 with 7.6 GB of memory available to it, using g++ 15.2
with \texttt{-O2}. Threads were not pinned to cores. For CPU-bound tests
the log is written to a tmpfs, where fsync costs nothing; the durability
tests write to an ext4 file system. That all results come from one machine
is a real limitation (Section~\ref{sec:limits}).

\paragraph{Data} Synthetic keys are uniform (1M and 10M keys, drawn from
$[0, 10n)$) and skewed (lognormal with $\sigma = 2$, scaled to $10^{15}$).
Real keys come from four SOSD datasets~\cite{marcus2020}: books, fb, wiki
and osm, each downsampled by keeping every $k$-th key to 10M keys, the way
SOSD itself shrinks its larger datasets, and checked against the checksums
SOSD publishes. wiki contains repeated timestamps, of which 2,417 were
dropped to make the keys unique; osm's keys were shifted right one bit so
they fit in a signed 64-bit integer. For SQL-level tests, keys too large
for a 48-bit primary key are shifted right until they fit, which keeps
their order and the shape of the distribution.

\paragraph{Procedure} Every structure's answers are checked before it is
timed. Within one program, each pass times every structure once, in turn,
and results are medians over seven passes; an earlier version that timed
structures one after another produced differences of more than 10\% for
the same structure depending on when it ran. The two builds of the engine
are separate programs, so each engine experiment runs three times per build,
alternating, and reports medians. On this machine the same measurement moves
by 10--20\% between runs, so differences smaller than that are not
meaningful. One script reruns every experiment and records the machine,
compiler and commit next to the raw output.

\section{Results}
\label{sec:results}

All numbers in this section come from one run of the benchmark script at
commit \texttt{f146a6d}, except Table~\ref{tab:spread}, which comes from a
later run of the insert benchmarks alone. Speedups are the B+-tree's time
divided by the learned index's, so values above 1 favor the learned
index. ``tlx'' is tlx's B+-tree with 512-byte nodes, its fastest setting
for lookups on this data among node sizes from 128 bytes to 2 KB.

\subsection{The index on its own}

Table~\ref{tab:index} times point lookups on the gapped array the engine
uses, freshly built on 90\% of each dataset, against tlx bulk-loaded on
the same keys and against binary search over a sorted array. The learned
index is 1.5--1.6x faster than tlx at 1M keys and 1.1--1.4x at 10M,
where neither fits in the processor's cache. On the real datasets the
range is wider: 1.30x faster on books, about even on fb, and slower on
osm. fb and osm are the two that need many more segments (their models
are five to twelve times larger per key than the other two real
datasets'), which matches their reputation as hard datasets for learned
indexes~\cite{marcus2020}.

\begin{table}[t]
\caption{Point lookups on the index alone (ns, medians of seven
interleaved passes) and the learned model's size.}
\label{tab:index}
\centering
\small
\setlength{\tabcolsep}{4pt}
\begin{tabular}{lrrrrr}
\toprule
Data & Learned & tlx & Binary & Speedup & Model (B/key) \\
\midrule
uniform 1M & 128.4 & 204.3 & 175.5 & 1.59 & 0.0065 \\
skewed 1M & 124.1 & 181.3 & 173.4 & 1.46 & 0.0074 \\
uniform 10M & 308.7 & 425.1 & 453.2 & 1.38 & 0.0064 \\
skewed 10M & 397.6 & 434.4 & 560.7 & 1.09 & 0.0066 \\
books 10M & 322.9 & 419.0 & 585.0 & 1.30 & 0.0060 \\
fb 10M & 393.8 & 410.1 & 565.9 & 1.04 & 0.0695 \\
osm 10M & 490.0 & 445.4 & 481.3 & 0.91 & 0.0723 \\
wiki 10M & 366.1 & 452.9 & 519.1 & 1.24 & 0.0147 \\
\bottomrule
\end{tabular}
\end{table}

After the remaining 10\% of the keys are inserted, the learned index is
1.3--2.0x faster than tlx on every dataset: tlx's leaves start out full
and split, while the gapped array had room set aside. The learned model
takes 0.006--0.07 bytes per key, against 0.25 for tlx's inner nodes (0.53
with its default 256-byte nodes), but the gapped array as a whole takes
11.4--11.5 bytes per key, against 8.6 for the packed tlx tree, because a
third of its slots are empty. Against the published PGM-index on 1M
keys~\cite{ferragina2020}, the pinned segmentation took 98.7 and 98.1~ns
on uniform and skewed data, and the PGM-index 112.7 and 122.2~ns; in an
earlier run the PGM-index took 97.5~ns on uniform data, so we read that
comparison as even on uniform data and in the learned index's favor on
skewed data.

\subsection{Reads inside the engine}

Table~\ref{tab:engine} compares the two builds of the whole engine on a
table of \texttt{(id INT, val INT)} rows. At 1M keys the store's own
lookup, which also copies the row out, is 1.6x faster on the learned
index, about as much as the index alone. Once the executor has planned
the query and decoded the row, the speedup is 1.12--1.18x, and once the
statement has to be parsed it is 1.05--1.08x: parsing, planning and
decoding cost more than the lookup. At 10M keys the store-level lookups
came out even in this run while the full queries were 15\% faster on
the learned index, an inconsistency we cannot explain and attribute to
noise. Lookups for ids that are not in the table are 1.8--2.3x faster on
the learned index on every dataset; we have not established why. Range
scans are even at every length, since decoding the rows dominates.

\begin{table}[t]
\caption{The whole engine on each index (ns per operation unless noted;
medians of three alternating runs of each build). Speedup is the
B+-tree's time over the learned index's.}
\label{tab:engine}
\centering
\small
\setlength{\tabcolsep}{3pt}
\begin{tabular}{p{4.0cm}rrr}
\toprule
& \multicolumn{3}{c}{Speedup} \\
Operation & uniform 1M & skewed 1M & uniform 10M \\
\midrule
store lookup (row copied out) & 1.56 & 1.60 & 0.97 \\
point query, parsed ahead of time & 1.12 & 1.18 & 1.15 \\
point query, parsed each time & 1.06 & 1.05 & 1.08 \\
point query, missing id & 2.30 & 2.15 & 2.31 \\
range scan, 100 rows & 0.98 & 1.01 & 0.99 \\
insert, random order & 0.93 & 0.94 & 0.99 \\
insert, appended & 0.52 & 0.51 & 0.33 \\
95\% reads, 5\% inserts following the data & 1.42 & 1.13 & 1.36 \\
50\% reads, 50\% inserts following the data & 1.06 & 0.89 & 1.11 \\
50\% reads, 50\% inserts spread evenly & 0.79 & 0.46 & 1.09 \\
durable insert, fsync every write & 1.00 & 1.03 & 1.26 \\
durable insert, fsync every 64 writes & 0.80 & 1.13 & 0.96 \\
open from a snapshot & 1.55 & 1.59 & 1.40 \\
\midrule
memory per row from a snapshot (B) & 89.2 / 94.3 & 89.2 / 94.3 & 89.2 / 94.3 \\
\bottomrule
\end{tabular}
\end{table}

\subsection{Writes}

Random-order inserts through SQL are even (0.93--0.99), as at the storage
level, where the gapped array takes 398 and 482~ns per insert on uniform
and skewed 1M keys against 326 and 385~ns for tlx. Appends past the end
of the key range, the pattern of an auto-increment primary key, are 2--3x
slower on the learned index. At the storage level the gap is larger, 12
to 15 times (466--577~ns against 38~ns for tlx), because a B+-tree
appends by touching only its rightmost leaf, while the gapped array keeps
shifting keys at its end and refreshing its last segment.

The mixed workloads separate two kinds of insert. When new ids are drawn
between existing ones, so that they follow the data's distribution, a
read-heavy mix (95\% reads) is 1.13--1.42x faster on the learned index,
which is consistent with tlx's lookups slowing down as its leaves split. When new
ids are drawn evenly from the whole key range, the learned index is up to
2.2x slower on skewed data. Table~\ref{tab:spread} shows why at the
storage level. Skewed data leaves a few segments that cover most of the
key range but hold few keys; evenly spread inserts pour into their small
regions, which fill at once and keep borrowing room from their
neighbors, so each insert re-lays far more of the array.

\begin{table}[t]
\caption{Inserting 10\% more keys at the storage level, either drawn
from the data or drawn evenly from the whole key range, in random order
(ns per insert, learned / tlx, and slots the gapped array re-lays per
insert). From run \texttt{2026-10-25\_1145} at commit \texttt{41d3f3c}.}
\label{tab:spread}
\centering
\small
\setlength{\tabcolsep}{4pt}
\begin{tabular}{lrrr}
\toprule
Data & From the data & Spread evenly & Slots re-laid \\
\midrule
uniform 1M & 389 / 292 & 385 / 322 & 49 $\to$ 48 \\
skewed 1M & 395 / 369 & 1,110 / 157 & 49 $\to$ 161 \\
books 10M & 570 / 546 & 986 / 572 & 47 $\to$ 77 \\
fb 10M & 456 / 516 & 1,452 / 201 & 3 $\to$ 129 \\
osm 10M & 445 / 561 & 1,277 / 289 & 5 $\to$ 85 \\
wiki 10M & 574 / 579 & 1,070 / 436 & 35 $\to$ 71 \\
\bottomrule
\end{tabular}
\end{table}

On uniform keys the two kinds of insert are the same thing, and the
numbers agree. Everywhere else, spreading the inserts makes the gapped
array 1.7--3.2x slower, raises the slots it re-lays per insert by up to
40 times, and forces refreshes that borrow from neighboring segments
(between 24 and 1,814 of them, against none when the inserts follow the
data). tlx stays about the same on books and gets faster on the others,
by up to 2.6x, presumably because the new keys land in a few leaves that
stay in cache.

Concentrating all inserts into a narrow key range (1.2--5\% of it) made
them 57--72\% slower than spreading them out (658--706 against 413~ns),
with no full rebuilds. This is a far weaker attack than published ones on
learned indexes~\cite{yang2024}; we did not attempt a targeted one.

\subsection{Durability}

With an fsync on every write, a durable insert takes 1.3--1.9~ms on
either index, and the index makes no measurable difference. Batching
helps as expected: on 1M uniform keys, fsyncing every 8, 64 and 512
writes brings the learned index's insert to 204, 38 and 7.6~$\mu$s (176,
30 and 7.7~$\mu$s for tlx), approaching the 2.6--2.8~$\mu$s an insert
costs with no fsync at all. At every batch size the two indexes are
within the run-to-run noise of each other.

\subsection{Startup and memory}

Opening the engine from a snapshot is 1.4--1.6x faster on the learned
index (140 against 217~ms at 1M rows, 1.7 against 2.4~s at 10M), and
from the log, which the engine folds into the snapshot's rows before a
single build, 1.3--1.8x faster. Laying out a sorted array is cheaper than
building a tree; the B+-tree wrapper also makes a second pass to move the
rows into the leaves, because tlx's bulk load copies its input. Memory
per row is about the same once both are built from a snapshot (89 against
94 bytes; most of a row is the row itself), and 26--30\% lower on the
learned index after a load by random-order inserts, which leaves tlx's
leaves partly empty.

\subsection{Real datasets}

On the four SOSD datasets (one run of each build at 10M keys, so noisier
than the medians above), a parsed point query was between 1.08x faster
and 1.01x slower on the learned index, appends were about 3x slower
(speedup 0.31--0.32), the read-heavy mix following the data was
1.21--1.62x faster, and inserts spread evenly over the key range ranged
from even (books) to 3.4x slower (fb).

\section{Lessons from Building It}
\label{sec:lessons}

\paragraph{A weak baseline hid quadratic inserts} For most of the
project's life, the insert benchmark compared the gapped array with
inserting into a plain sorted vector, which shifts every later key, and
reported a twelvefold advantage. Compared later against tlx, inserts were
45 to 1,200 times slower than the B+-tree, and the gap grew with the array.
The cause was a safety net: a full rebuild every 512 inserts regardless of
the array's size, which made the average insert linear in the number of
keys. It existed because the drift accounting ignored shifts that crossed
into a neighboring segment. Counting those, adding windowed rebalancing and
growth at the end, and capping segment length removed every full rebuild
from the benchmarks.

\paragraph{A third of every lookup was a mispredicted branch} Finding a
key's segment used \texttt{std::upper\_bound}, whose branch goes either way
with equal probability for random keys. Rewriting the binary search to
select each half with a conditional move cut it from about 48 to about
9~ns, a third of the lookup time at 1M keys, and changed the comparison with
the PGM-index from a clear loss on uniform data to a tie.

\paragraph{An invariant's test must fail without the mechanism} A test of
the $2\epsilon$ drift bound passed with the neighbor counting removed. The
nearest-gap search alternates sides, which limits how far shifts push into a
neighbor unless the neighbor has a long run of empty slots at its edge. Only
after the test deleted keys next to a segment boundary did it fail without
the counting (worst error $2.47\epsilon$) and pass with it.

\paragraph{Differential testing found what unit tests did not} Running
random statements against both this engine and SQLite, with automatic
shrinking of failing cases, found a \texttt{LIMIT} applied before
\texttt{COUNT(*)}, secondary-index entries left behind by inserts and
updates that replaced rows, and out-of-range values silently truncated into
other keys. The same test now runs on both builds of the engine, and with a
tiny error bound ($\epsilon = 4$) that makes the rebalancing code run
constantly.

\paragraph{The best number came from a broken dataset} An early skewed
dataset collapsed into a run of consecutive integers, which a single line
fits exactly, and produced the project's fastest lookups. Looking at the
gaps between keys, which nobody had done, would have caught it immediately.

\section{Limitations}
\label{sec:limits}

All results come from one laptop under WSL2, with threads not pinned to
cores. The engine keeps everything in memory and reads its whole snapshot
on startup, unlike disk-based engines that load pages on demand. It has one
client, no multi-statement transactions and no joins. Only one B+-tree
(tlx) and one learned index design are compared; an updatable learned index
such as ALEX inside the same engine would be a natural next comparison. The
real datasets are downsampled to 10M keys, and SQL-level tests shift large
keys to fit 48 bits. Because the two builds are separate programs, their
comparison relies on alternating runs rather than interleaved passes.

\section{Conclusion}

Swapping only the index under one engine shows how much of a learned
index's advantage reaches the user of a small embedded database: a 1.5x
faster lookup at the storage layer becomes a 5--8\% faster query once it
has to be parsed, while the learned index's weaknesses on appends and on
inserts that depart from the data's distribution carry through almost in
full. For an engine like this one, a learned index is not a clear win
over a tuned B+-tree. It trades slightly faster lookups, faster startup
and better behavior after random inserts for much slower appends and a
sensitivity to where inserts land. The obvious next steps are an
updatable learned index such as ALEX in the same engine, a disk-resident
version, and measurements on more than one machine.

\section*{Artifact}

The code, the scripts that produce every number in this paper, and the raw
output of the runs reported here are at
\url{https://github.com/asood8/Learned-Index}; the runs reported here
were made at commits \texttt{f146a6d} and \texttt{41d3f3c}, and their
output is in \texttt{results/runs/2026-10-25\_1112} and
\texttt{results/runs/2026-10-25\_1145}. \tbd{Zenodo DOI of the tagged
release.}

\bibliographystyle{IEEEtran}
\bibliography{refs}

\end{document}
